\chapter*{TÓM TẮT}
\addcontentsline{toc}{chapter}{TÓM TẮT}
\fontsize{13}{15}\selectfont

\textbf{Tóm tắt:} Khóa luận này trình bày các kết quả nghiên cứu và quy trình thiết kế một hệ thống giám sát và điều khiển chất lượng không khí trong nhà (Indoor Air Quality - IAQ) sử dụng cây trồng thủy sinh, được thiết kế chuyên biệt cho không gian kín như căn hộ chung cư. Nghiên cứu có sự kết hợp  giữa nền tảng Internet vạn vật (IoT) và các giải pháp thực sinh học, cụ thể là ứng dụng cây lan ý (\textit{Spathiphyllum}). Khác với các hướng tiếp cận truyền thống vốn chỉ xem cây như một màng lọc khí thụ động,nghiên cứu này coi thực vật là một hệ thống sinh học tự động có khả năng tương tác. Khả năng này chỉ được phát huy tối đa khi môi trường sinh trưởng được duy trì ở trạng thái lý tưởng. Kế thừa các công bố khoa học về ứng dụng thủy canh trong kiểm soát IAQ, lan ý đã được chứng minh là loài thực vật sở hữu năng lực vượt trội trong việc đồng hóa CO$_2$ và hấp thụ các hợp chất hữu cơ dễ bay hơi (VOC) tích tụ trong không gian kín \cite{hydroponic_iaq_survey,nasa_clean_air}. Phát triển từ đó, mục tiêu cốt lõi của đề tài là  tạo một hệ sinh thái khép kín: vừa giám sát, cảnh báo tự động, vừa điều phối các điều kiện vi khí hậu nhằm tối ưu hóa sự phát triển của thực vật và bảo vệ sức khỏe con người.

Về mặt kiến trúc hệ thống, nghiên cứu triển khai mô hình phân lớp chuẩn mực bao gồm: sensor node, gateway, cloud và application. Trong đó, node sử dụng vi điều khiển Arduino Pro Mini để đo lường  các tham số môi trường và sinh trưởng như: chỉ số tổng chất rắn hòa tan (TDS), độ pH, nhiệt độ và độ ẩm không khí, nồng độ CO$_2$ và cường độ ánh sáng. Nhằm khắc phục triệt để nhược điểm về vùng phủ sóng của mạng Wi-Fi truyền thống trong cấu trúc căn hộ nhiều vật cản, hệ thống đã  ứng dụng công nghệ truyền thông công suất thấp, khoảng cách xa LoRa để kết nối node với gateway . Tại lớp trung gian này, gateway không chỉ đóng vai trò chuyển tiếp mà còn thực thi tiền xử lý, kiểm định tính toàn vẹn của dữ liệu trước khi đồng bộ lên nền tảng đám mây ThingsBoard. Thingsboard cũng cung cấp giao thức giao tiếp ứng dụng (API) cho phần mềm di động được phát triển trên Flutter. Từ đó, hệ thống trao quyền cho người dùng khả năng giám sát theo thời gian thực và trực tiếp can thiệp vào các thiết bị ngoại vi (quạt thông gió, đèn chiếu sáng, điều hòa không khí) thông qua  rơ-le và sóng hồng ngoại.

Những đóng góp khoa học và thực tiễn của nghiên cứu tập trung vào việc giải quyết bài toán tích hợp hệ thống: từ thiết kế phần cứng, chuẩn hóa giao thức truyền thông, kiến trúc dữ liệu trên ThingsBoard cho đến việc thiết kế application . Để đảm bảo sự sinh trưởng bền vững của hệ thủy canh trong không gian hẹp, hệ thống sử dụng các ngưỡng giới hạn nghiêm ngặt về pH (6.0 – 6.5) và độ dẫn điện (EC) ở mức xấp xỉ 1.2 dS/m(768ppm), dựa trên các tiêu chuẩn nông nghiệp hiện hành \cite{peace_lily_hydro_brightlane,peace_lily_ec_ppf}. Các chuỗi thực nghiệm đánh giá đã khẳng định độ tin cậy và khả năng đáp ứng yêu cầu của hệ thống trước những biến động bất lợi của môi trường, điển hình như sự gia tăng đột biến của CO$_2$, thiếu hụt ánh sáng hay sự mất cân bằng chất dinh dưỡng. Kết quả này không chỉ mang tính ứng dụng cao mà còn tạo tiền đề vững chắc cho các nghiên cứu tiếp theo về ứng dụng trí tuệ nhân tạo trong kiểm soát không gian sống.

\textbf{\textit{Từ khóa: }} \textit{Chất lượng không khí trong nhà (IAQ), Internet vạn vật (IoT), lan ý, Spathiphyllum, hệ thủy canh, Arduino Pro Mini, ESP32, truyền thông LoRa, ThingsBoard, Flutter, nồng độ CO2, điều khiển thích nghi.}
